\section{Introduction}


Digital marketing presents a vast optimization landscape. A single Facebook advertising campaign involves dozens of parameters: headline text, body copy, images, call-to-action buttons, audience targeting (demographics, interests, behaviors), placement (feed, stories, right column), and bidding strategy. The combinatorial explosion of these parameters renders exhaustive search infeasible.

Current practice relies on A/B testing: marketers hypothesize improvements, create variants, and measure performance. This approach suffers from three limitations:

\begin{enumerate}
\item \textbf{Limited exploration}: Human creativity bounds the variants considered.
\item \textbf{Sequential evaluation}: Testing one hypothesis at a time delays learning.
\item \textbf{Local optima}: Incremental improvements miss global structure.
\end{enumerate}

Earle addresses these limitations by applying genetic algorithms to marketing optimization. We model campaigns as genomes and apply evolutionary operators—selection, crossover, mutation—to explore the parameter space efficiently.

Our contributions are:
\begin{enumerate}
\item A formal model of marketing campaigns as structured genomes amenable to genetic optimization.
\item Domain-specific genetic operators for text, image, and audience evolution.
\item A multi-objective fitness function balancing conversion, cost, and brand consistency.
\item Empirical evaluation demonstrating significant improvements over human-designed campaigns and alternative optimization methods.
\end{enumerate}

